\section{The Dual-Constraint Framework}
\label{sec:framework}

\subsection{Setup}

Let $G = (V, E, w)$ be a state's census-tract adjacency graph.
Each vertex $v \in V$ carries:
\begin{itemize}
  \item $p_v > 0$: total census population
  \item $d_v \geq 0$: two-party Democratic votes (2020 presidential)
\end{itemize}

Let $P = \sum_v p_v$, $D = \sum_v d_v$, $R = P - D$ (two-party totals).
Write $d = D/P$ for the statewide Democratic vote share and $r = 1 - d$.

Let $k$ be the congressional district count.
A \emph{proportional seat allocation} assigns $k_D = \mathrm{HH}(d, r, k)$
seats to Democrats and $k_R = k - k_D$ to Republicans, where $\mathrm{HH}$
is the Huntington-Hill rounding.

\subsection{Proportionality Gap}

\begin{definition}[Proportionality gap]
\label{def:pgap}
Let $d = D/(D+R)$ be the statewide Democratic two-party vote share, $k$ the
total congressional seat count, and $s_D$ the number of Democratic-won seats
in the final plan.
The \emph{proportionality gap} is
\[
  \Delta = \frac{s_D}{k} - d.
\]
$\Delta > 0$ indicates Democratic over-representation; $\Delta < 0$ indicates
Republican over-representation.
The proportionally fair outcome is $\Delta = 0$, achieved when each party's
seat share equals its vote share.
\end{definition}

\noindent Note that $\Delta$ differs from the efficiency gap \citep{mcghee2014measuring}:
the efficiency gap is based on wasted votes ($\text{EG} = (W_D - W_R)/P$),
whereas $\Delta$ directly compares seat share to vote share.
The two measures agree in sign for strongly disproportionate outcomes but
diverge for near-proportional plans; Section~\ref{sec:legal} discusses the
implications for legal standards.

\subsection{The ncon=2 Bisection}

We extend the minimum-edge-cut bisection to METIS's dual-constraint mode with
vertex weights $\mathbf{w}_v = (p_v,\, d_v)$.
A bisection partitions $V$ into $(S,\, V \setminus S)$ subject to:
\begin{align}
  \frac{\sum_{v \in S} p_v}{P} &\in \left[\frac{k_D}{k} - \varepsilon_p,\; \frac{k_D}{k} + \varepsilon_p\right] \label{eq:pop-balance}\\
  \frac{\sum_{v \in S} d_v}{D} &\in \left[\tau_D - \varepsilon_d,\; \tau_D + \varepsilon_d\right] \label{eq:dem-balance}
\end{align}
where $\varepsilon_p = 0.001$ (0.1\,\% population tolerance), $\varepsilon_d$ is
the partisan tolerance parameter, and $\tau_D$ is the target Democratic fraction
for the left (D-bloc) partition.

The target is implemented via METIS's $\mathtt{tpwgts}$ array:
\[
  \mathtt{tpwgts} = [\underbrace{k_D/k}_{\text{left pop}},\; \underbrace{\tau_D}_{\text{left D}},\;
                     \underbrace{k_R/k}_{\text{right pop}},\; \underbrace{1-\tau_D}_{\text{right D}}]
\]

\subsection{The Proportionality Formula}
\label{sec:formula}

The key design question is: what value of $\tau_D$ guarantees that the D-bloc
produces $k_D$ Democratic-winning districts and the R-bloc produces $k_R$
Republican-winning districts?

\begin{definition}[Proportional target]
A target $\tau_D^*$ is \emph{proportionally optimal} if:
\begin{enumerate}
  \item The right (R-bloc) has exactly 50\,\% Democratic concentration:
    $\frac{(1-\tau_D) D}{(k_R/k) P} = 0.5$
  \item No more Democratic votes are allocated to the R-bloc than this minimum.
\end{enumerate}
\end{definition}

Solving condition~(1):
\[
  (1 - \tau_D^*) \cdot d = 0.5 \cdot \frac{k_R}{k}
  \quad \Longrightarrow \quad
  \tau_D^* = 1 - \frac{k_R}{2dk}
\]

The full $\mathtt{tpwgts}$ vector for proportional bisection is therefore:
\begin{equation}
  \label{eq:tpwgts-formula}
  \boxed{\mathtt{tpwgts} = \left[\frac{k_D}{k},\; 1 - \frac{k_R}{2dk},\;
                                  \frac{k_R}{k},\; \frac{k_R}{2dk}\right]}
\end{equation}

\noindent\textbf{Example (60/40 state, $k=3$, $k_D=2$, $k_R=1$):}
\[
  \mathtt{tpwgts} = \left[\frac{2}{3},\; 1 - \frac{1}{3.6},\; \frac{1}{3},\; \frac{1}{3.6}\right]
                  = [0.667,\; 0.722,\; 0.333,\; 0.278]
\]

D-bloc Democratic concentration: $0.722 \times 0.6 / (2/3) = 65\,\%$ D.
R-bloc Democratic concentration: $0.278 \times 0.6 / (1/3) = 50\,\%$ D. \checkmark

\subsection{The Partisan Tolerance Parameter}

METIS's $\mathtt{ubvec}[1]$ controls how strictly the Democratic-vote constraint
is enforced:
\begin{itemize}
  \item $\mathtt{ubvec}[1] = 1.0$: hard constraint (exact proportionality required)
  \item $\mathtt{ubvec}[1] = 1.05$: 5\,\% tolerance around target
  \item $\mathtt{ubvec}[1] \to \infty$: constraint disabled (reverts to B.9/B.8)
\end{itemize}

The \emph{compromise parameter} $\mathtt{ubvec}[1]$ traces a family of algorithms
from pure geographic neutrality ($\infty$) to strict partisan proportionality (1.0).
Section~\ref{sec:eval} characterises the tradeoff curve for each state.

\subsection{Connection to AreaSection}

AreaSection (B.9) uses the identical $\mathtt{ncon}=2$ infrastructure with
vertex weights $[\text{population},\, \text{area}]$ and
$\mathtt{tpwgts} = [k_L/k,\; 0.5,\; k_R/k,\; 0.5]$ (targeting 50\,\%
land area on each side).

B.12 substitutes $D_{\text{votes}}$ for land area, and replaces the fixed
$0.5$ area target with the derived $\tau_D^*$ that achieves proportionality.

The algorithmic infrastructure is identical.
The political implications differ entirely:
land-area balance has no partisan interpretation; Democratic-vote balance
is explicitly partisan.
